\section{Perturbation theory}\label{sec:perturbation}

The $n$-point
correlation functions of the field $B^a_{\mu}(t,x)$
can be computed in perturbation theory by expanding the
field in powers of the fundamental field, as explained sect.~\ref{sec:iterative},
and by calculating the correlation functions of the latter
in the $\SUn$ gauge theory as usual.
We now show that the correlation functions can be
directly obtained from
a set of Feynman rules in $D+1$ dimensions.

\subsection{Gauge fixing}\label{ssec:gaugefixing}

The flow equation \eqref{eq:2.4}
is invariant under the infinitesimal transformation
\begin{equation}
  \dbrs B_{\mu}=D_{\mu}\omega,
  \label{eq:3.1}
\end{equation}
provided the time-dependence of $\omega(t,x)\in\sun$
is such that
\begin{equation}
  \partial_t\omega=\alpha_0D_{\mu}\partial_{\mu}\omega.
  \label{eq:3.2}
\end{equation}
Since the initial value of $\omega$ is unconstrained, these transformations
generate the full gauge group at flow time $t=0$ and thus extend the
gauge symmetry of the $\SUn$ gauge theory to all flow times.

The symmetry can be fixed as usual by including the
gauge-fixing term
\begin{equation}
  \Sgf=-{\lambda_0\over g_0^2}\int\rmd^Dx\,
  \tr\{\partial_{\mu}A_{\mu}(x)\partial_{\nu}A_{\nu}(x)\}
  \label{eq:3.3}
\end{equation}
and the associated ghost action
\begin{equation}
  \Sgh=-{2\over g_0^2}\int\rmd^Dx\,
  \tr\{\partial_{\mu}\cbar(x)D_{\mu}c(x)\}
  \label{eq:3.4}
\end{equation}
in the total action of the theory, where $c$ and $\cbar$ are the
Faddeev--Popov ghost fields.
As far as the $n$-point
correlation functions of the fundamental gauge field
and the ghosts are concerned, the Feynman rules are then the standard
ones. Evidently, since the transformation \eqref{eq:3.1} is
an infinitesimal gauge variation, the expectation values of
gauge-invariant expressions in the field generated by the flow
are independent of the gauge parameter $\lambda_0$.

\subsection{Gauge-field propagator}

The way in which the flow-line diagrams combine with the
Feynman diagrams of the underlying theory is best
explained by considering the two-point function
of the time-dependent gauge field.
To leading order in the gauge coupling,
the flow-line diagram
\begin{equation}
  \raisebox{-0.5cm}{\includegraphics[width=3.5cm]{images/diagram2}}
  \label{eq:3.5}
\end{equation}
is the only one that contributes to the correlation function.
The contraction of the gauge fields at the one-point
vertices then shows that
\begin{align}
  \langle \tilde{B}^a_{\mu}(t,p)\tilde{B}^b_{\nu}(s,q)\rangle=
  (2\pi)^D\delta(p+q)\delta^{ab}g_0^2\Dtilde_{t+s}(p)_{\mu\nu}+\rmO(g_0^4),
  \label{eq:3.6}\\[2pt]
  \Dtilde_t(p)_{\mu\nu}=
  {1\over(p^2)^2}
  \bigl\{
  (\delta_{\mu\nu}p^2-p_{\mu}p_{\nu})\rme^{-tp^2}+
  \lambda_0^{-1}p_{\mu}p_{\nu}\rme^{-\alpha_0tp^2}\bigr\}.
  \label{eq:3.7}
\end{align}
This formula includes the mixed propagator
\begin{equation}
  \langle \tilde{A}^a_{\mu}(p)\tilde{B}^b_{\nu}(s,q)\rangle=
  \langle \tilde{B}^a_{\mu}(0,p)\tilde{B}^b_{\nu}(s,q)\rangle
  \label{eq:3.8}
\end{equation}
as well as the two-point function of the gauge field
at flow time zero.

Since all three propagators are given by the same analytic expression,
the same graphical symbol
\begin{equation}
  \raisebox{-0.55cm}{\includegraphics[width=2.5cm]{images/prop2}}%
  =\;\delta^{ab}g_0^2\Dtilde_{t+s}(p)_{\mu\nu}
  \label{eq:3.9}
\end{equation}
may be used for them. Note that the contraction of the one-point vertices
always has the effect of converting the terminal flow lines
to gauge-field lines.
If one starts from the flow-line diagram
\begin{equation}
  \raisebox{-0.6cm}{\includegraphics[width=4.5cm]{images/diagram4}}
  \label{eq:3.10}
\end{equation}
instead of the diagram \eqref{eq:3.5}, for example, the
leading-order contribution is obtained by substituting
the tree-level diagram for the correlation function
of the gauge fields at the three one-point vertices.
This leads to a diagram
\begin{equation}
  \raisebox{-0.5cm}{\includegraphics[width=3.6cm]{images/diagram5}}
  \label{eq:3.11}
\end{equation}
that has an ordinary three-point vertex (filled circle)
and a flow vertex. Two of~the lines
attached to the latter are gauge-field lines instead
of flow lines, but the expression for the vertex is the same as before.

\subsection{Feynman rules in $D+1$ dimensions}

The flow time will now be interpreted as an additional
space-time coordinate. Since only non-negative times are considered,
the $D+1$ dimensional space is a half-space with
a $D$ dimensional boundary at flow time zero.
The $\SUn$ gauge theory
lives at the boundary, while the field generated by the
gradient flow extends to the extra dimension.

From this point of view, the ordinary and the flow vertices represent
boundary and bulk interaction terms, respectively, while
the propagation of the fields in $D+1$ dimensions is described
by the gauge-field propagator \eqref{eq:3.9} and the flow propagator \eqref{eq:2.16}.
The ghost-field propagator
\begin{equation}
  \langle\ctilde^a(p)\cbartilde^b(q)\rangle=
  (2\pi)^D\delta(p+q)\delta^{ab}g_0^2\Dtilde(p)+\rmO(g_0^4),
  \qquad
  \Dtilde(p)={1\over p^2},
  \label{eq:3.12}
\end{equation}
on the other hand, is defined at flow time zero only.
Graphically it is represented through a dotted line,
\begin{equation}
  \raisebox{-0.55cm}{\includegraphics[width=1.8cm]{images/prop3}}%
  =\;\delta^{ab}g_0^2\Dtilde(p),
  \label{eq:3.13}
\end{equation}
where the arrow is drawn in the direction
from $\cbar$ to $c$.

\begin{figure}[t]
\centerline{\includegraphics[width=10.0cm]{images/diagram6}}
\caption{Examples of diagrams contributing to the two- and three-point
functions of the time-dependent gauge field.
All diagrams are built from flow propagators
(directed solid lines), gauge-field propagators (wiggly lines),
ghost-field propagators (directed dotted lines), flow vertices (open circles)
and ordinary vertices at flow time zero (filled circles). The little
square at the end of the external lines indicates that
they are not amputated.
}
\label{fig:examples}
\end{figure}

It should be quite clear at this point that the $n$-point functions
of the field generated by the gradient flow are given by
Feynman diagrams in $D+1$ dimensions with the graphical
elements listed in the caption of fig.~\ref{fig:examples}.
Some special features of the Feynman rules are worth pointing out:

\begin{enumerate}[label=(\alph*),itemsep=1.0ex,leftmargin=*]
\item Each flow vertex has exactly one outward-directed flow line corresponding
to the first momentum-index combination in eq.~\eqref{eq:2.13}. The other lines
attached to these vertices may be gauge-field lines or ingoing flow lines.

\item Flow lines must start at a
flow vertex and are either external or end at another flow vertex.
The gauge-field lines, on the other hand, can start and end at
both the flow vertices and the ordinary vertices.

\item The flow vertices are inserted at some flow time which is integrated
from zero to infinity. All other vertices are at flow time zero.
The flow times on which the propagators depend are the ones at
the endpoints of the corresponding lines.

\item Diagrams with closed flow-line loops are set to zero.
This rule derives from the fact that flow-line diagrams are tree
diagrams and that the contraction of the gauge fields at the one-point
vertices never leads to new flow lines.
\end{enumerate}
